% !TeX program = pdflatex
% ==========================================================================
%  HPDC-MaOEA -- two-column manuscript draft (Elsevier elsarticle, 5p layout)
%
%  The manuscript is organised around two contributions: hybrid PBI quality
%  grouping and dual-mode candidate selection. Established components are
%  cited where they enter the framework.
%
%  Layout notes:
%    * elsarticle option "5p" is Elsevier's final two-column journal layout
%      (textwidth 522pt, columnsep 18pt => 252pt columns) and it loads
%      fleqn.clo, so displayed equations are flush left.
%    * Wide displays are broken with split/multline and wide floats use the
%      starred environments (table*, algorithm*) to span both columns.
%      Line breaking and float placement are adapted to the column width; the
%      wording and the mathematics are not affected by it.
%    * For the double-spaced single-column referee version, swap the class
%      options for [preprint,review,12pt].
% ==========================================================================
\documentclass[final,5p,times,twocolumn]{elsarticle}

\usepackage{amsmath,amssymb,amsthm}
\usepackage{bm}
\usepackage{booktabs}
\usepackage{algorithm}
\usepackage{algpseudocode}
\usepackage{enumitem}
\usepackage{xcolor}
\usepackage[hidelinks]{hyperref}

% graphicx is already loaded by elsarticle.cls

\newtheorem{proposition}{Proposition}
\newtheorem{remark}{Remark}

% Narrow columns: allow slightly looser lines before an overfull box.
\tolerance=1200
\emergencystretch=1em

% Pseudo-code is set one size down so that the wider \State lines fit a column.
\usepackage{etoolbox}
\AtBeginEnvironment{algorithmic}{\footnotesize}
\algrenewcommand{\algorithmicrequire}{\textbf{Input:}}
\algrenewcommand{\algorithmicensure}{\textbf{Output:}}

\newcommand{\vf}{\mathbf{f}}
\newcommand{\vx}{\mathbf{x}}
\newcommand{\vv}{\mathbf{v}}
\newcommand{\vw}{\mathbf{w}}
\newcommand{\vr}{\mathbf{r}}
\newcommand{\vz}{\mathbf{z}}
\newcommand{\TODO}[1]{\textcolor{red}{[TODO: #1]}}

%% Uncomment and fill in for a preprint that names its target journal.
%% \journal{Information Sciences}

\begin{document}

\begin{frontmatter}

\title{Hybrid PBI-Based Quality Grouping and Dual-Mode Candidate Selection
for Expensive Many-Objective Optimization}

%% Author block. Replace the placeholders and add \corref/\ead as needed:
%% \author[inst1]{Given Name Surname\corref{cor1}}
%% \ead{name@example.edu}
%% \cortext[cor1]{Corresponding author.}
\author{\TODO{author list}}

\affiliation{organization={Department, University},
            addressline={Street},
            city={City},
            postcode={000000},
            state={Province},
            country={Country}}

\begin{abstract}
\TODO{Use the challenge--contribution structure. State the two concrete gaps:
binary grouping alone provides no within-label ordering for positive-group selection,
and a fixed relation-score rule lets the score distribution determine how many
candidates are evaluated. Introduce hybrid PBI-based quality grouping and
dual-mode candidate selection as two coordinated design components, addressing
training-group construction and expensive-evaluation selection. Lead the results
with the strongest verified algorithm-level IGD comparison in the tested
many-objective settings. Then give one concise mechanism finding about group
composition or candidate-batch utility. Keep detailed overlap statistics and
componentwise tests in the experimental section.}
\end{abstract}

\begin{keyword}
expensive many-objective optimization \sep surrogate-assisted evolutionary
algorithm \sep relation learning \sep penalty-based boundary intersection \sep
candidate selection
\end{keyword}

\end{frontmatter}

\section{Introduction}
\label{sec:intro}

\TODO{Use a four-move opening: expensive many-objective relation learning; the
loss of within-label ordering when forming the positive group and the dependence of
fixed-cutoff batch size on the relation-score distribution; the two proposed
mechanisms; and algorithm-level performance followed by evidence explaining
group construction and candidate selection. Present the two mechanisms as
complementary design roles in one algorithm. Do not lead with selection-overlap
statistics or imply that the relation network learns a within-group ordering.}

The contributions are twofold. First, a hybrid PBI-based quality grouping
method combines a continuous PBI score with the binary representative-based
label used in REMO. The continuous score distinguishes solutions assigned the
same binary label when the positive training group is constructed. Second, a
dual-mode candidate selection method reduces the reliance of expensive
evaluations on a single relation score. Relation-guided candidate generation is
retained, while indicator-based reranking and exploration-oriented selection
provide two alternative criteria for selecting the final evaluation batch.

\section{Related Work}
\label{sec:related}

\subsection{Surrogate-Assisted Expensive Multi-/Many-Objective Optimization}
\label{sec:related:saea}
\TODO{Kriging/RBF-based SAEAs, their scalability limit in $M$, and why
regression on every objective becomes unattractive when $M$ grows.}

\subsection{Relation Learning for Expensive Optimization}
\label{sec:related:relation}
\TODO{REMO and the pairwise-comparison SAEA line; compare how each method forms
training groups or pairwise labels, then position the present work as a hybrid
grouping signal that combines continuous quality ranking with a binary
representative-based group label.}

\subsection{Quality Assessment and Candidate Selection}
\label{sec:related:selection}
\TODO{Indicator-based assessment (SDE, $L_p$ shape estimation, PIEA), radial
representative selection (RSEA), and infill/candidate selection rules. End with
the gap that motivates Section~\ref{sec:method}.}

\section{The Proposed HPDC-MaOEA}
\label{sec:method}

Relation-based surrogate-assisted evolutionary algorithms replace multi-output
regression by a classifier that predicts the relative quality of a pair of
solutions~\cite{ref:remo}. Two interfaces then decide what such an algorithm
can learn and how much its predictions cost: the rule that turns the evaluated
population into supervision labels, and the rule that turns the surrogate
scores of a large candidate pool into a small batch of expensive evaluations.
HPDC-MaOEA (hybrid-PBI dual-mode-candidate many-objective evolutionary
algorithm) contributes one mechanism at each interface. \emph{Hybrid PBI-based
quality grouping} (Section~\ref{sec:method:hpc}) combines a continuous PBI
score with a binary representative-based label. The score can distinguish
solutions assigned the same binary label when the positive training group is
constructed. \emph{Dual-mode candidate selection}
(Section~\ref{sec:method:candidate}) retains relation-guided candidate
generation but uses either indicator-based reranking or exploration-oriented
selection to choose the final expensive-evaluation batch. Established relation
learning, indicator assessment and environmental selection modules provide the
host algorithm; the two proposed mechanisms change how its training groups are
formed and how its surrogate scores are used.

\subsection{Overall Framework}
\label{sec:method:framework}

We consider the expensive many-objective problem
\begin{equation}
\min_{\vx\in\Omega}\;
\vf(\vx)=\bigl[f_1(\vx),\ldots,f_M(\vx)\bigr],
\qquad
\Omega\subseteq\mathbb{R}^{D},
\label{eq:problem}
\end{equation}
in which every evaluation of $\vf$ is expensive, so the number of real function
evaluations is limited to a budget $FE_{\max}$. Let $FE$ be the number of
evaluations consumed so far and $t=FE/FE_{\max}\in[0,1]$ the evaluation
progress, which both proposed mechanisms use as their only schedule variable.

HPDC-MaOEA evaluates a Latin hypercube design of $N_{\mathrm{init}}$ solutions,
stores them in an archive $\mathcal{A}_{\mathrm{rc}}$, and then repeats five
steps until the budget is exhausted. (i)~The current population $\mathcal{P}$
of size $N$ is divided into a positive group and a non-positive group by
hybrid PBI grouping, which also returns the set $\mathcal{R}$ of
representative solutions reused later as an additional mating pool.
(ii)~Ordered pairs are generated from the two groups and a three-class relation
network is trained, together with a held-out pairwise error rate $e_r$.
(iii)~An indicator value is computed for every evaluated solution and an
RBF-SVR model is fitted from decision vectors to those values.
(iv)~A relation-guided inner search accumulates a candidate pool, one of the
two selection criteria is drawn, and a batch of at most $n_{\max}$ candidates is
returned. (v)~The batch is evaluated, appended to
$\mathcal{A}_{\mathrm{rc}}$, and the next population is obtained by the
radial-grid environmental selection of RSEA~\cite{ref:rsea} over the whole
archive.

Algorithm~\ref{alg:framework} states the procedure. All expensive evaluations
after initialisation are consumed in step~(v), so both proposed mechanisms operate on already
evaluated data and surrogate predictions only, and neither changes the total
number of real evaluations. The two interfaces stay separate: grouping
decides \emph{what the classifier is asked to learn}, candidate selection
decides \emph{how its scores are spent}.

%% \begin{figure*}[t]
%%   \centering
%%   \includegraphics[width=\textwidth]{fig_framework}
%%   \caption{Overall framework of HPDC-MaOEA. \label{fig:framework}}
%% \end{figure*}
\TODO{Export the overall framework diagram as \texttt{fig\_framework.pdf} and
uncomment the figure* block above; it should show the outer loop of
Algorithm~\ref{alg:framework} with the two contributed blocks highlighted.}

\begin{algorithm}[t]
\caption{HPDC-MaOEA}
\label{alg:framework}
\begin{algorithmic}[1]
\Require problem, budget $FE_{\max}$, population size $N$
\Ensure archive $\mathcal{A}_{\mathrm{rc}}$
\State Sample $N_{\mathrm{init}}$ solutions by Latin hypercube design and
       evaluate them
\State $\mathcal{P}\gets$ initial population;
       $\mathcal{A}_{\mathrm{rc}}\gets\mathcal{P}$
\While{$FE<FE_{\max}$}
  \State $t\gets FE/FE_{\max}$
  \State $(\mathcal{C}_1,\mathcal{C}_2,\mathcal{R})\gets$
         \Call{HybridGroup}{$\mathcal{P},t$}
         \Comment{Section~\ref{sec:method:hpc}}
  \State Build ordered relation pairs from $\mathcal{C}_1,\mathcal{C}_2$ by
         \eqref{eq:pairlabel}
  \If{no pair can be formed}
    \State $\mathcal{P}\gets$
           \Call{EnvSelect}{$\mathcal{A}_{\mathrm{rc}},N$};
           \textbf{continue}
  \EndIf
  \State Train the relation surrogate and record its held-out error $e_r$
         \Comment{Section~\ref{sec:method:relation}}
  \State Train the indicator surrogate $\widehat I$ on the evaluated solutions
         \Comment{Section~\ref{sec:method:candidate:ind}}
  \State Draw the selection mode $m$ by \eqref{eq:mode}
  \State $\mathcal{S}\gets$
         \Call{DualModeSelection}{$\mathcal{P},\mathcal{R},m,t,e_r$}
         \Comment{Algorithm~\ref{alg:candidate}}
  \If{$\mathcal{S}=\emptyset$}
    \State $\mathcal{S}\gets$ at most $n_{\min}$ offspring generated from
           $\mathcal{P}\cup\mathcal{R}$
  \EndIf
  \State Truncate $\mathcal{S}$ by the remaining budget, evaluate it, and
         append it to $\mathcal{A}_{\mathrm{rc}}$
  \State $\mathcal{P}\gets$ \Call{EnvSelect}{$\mathcal{A}_{\mathrm{rc}},N$}
\EndWhile
\end{algorithmic}
\end{algorithm}

\subsection{Hybrid PBI-Based Quality Grouping}
\label{sec:method:hpc}

Relation learning first needs to divide the evaluated population into groups
from which pairwise labels can be generated. The binary PBI partition of
REMO~\cite{ref:remo} provides a clear group boundary, but all solutions assigned
the same label receive the same group-level value. It therefore cannot
distinguish two solutions of different PBI quality that lie on the same side of
the boundary. Hybrid PBI-based grouping combines a continuous PBI quality score
with the binary representative-based label. The continuous score provides
additional ranking information among solutions sharing a binary label, while
the label preserves a clear population-level partition. A stage-dependent
combination of these two signals determines the positive group used to construct
relation pairs.

Throughout this subsection $\mathcal{P}=\{(\vx_i,\vf_i)\}_{i=1}^{N}$ is the
current population and $\vz^{*}=\min_i\vf_i$ the component-wise population
minimum, taken as the ideal point.

\subsubsection{Continuous PBI Quality Score}
\label{sec:method:hpc:cont}

The continuous score assesses each solution by a PBI value measured along a
direction set $\mathcal{V}=\{\vv_j\}$. For low-dimensional or small-population
cases, $\mathcal{V}$ consists of uniform reference vectors normalised to unit
length. In the many-objective regime, the directions adapt to the current
nondominated distribution and are formed from its radial unit vectors,
\begin{equation}
\mathcal{V}=\Bigl\{\vf^{(j)}\big/\bigl\|\vf^{(j)}\bigr\|_2
\;\Bigm|\;j=1,\ldots,n_{\mathrm{ND}}\Bigr\},
\label{eq:direction}
\end{equation}
which are used directly as the adaptive direction set. Equation
\eqref{eq:direction} makes both the number and positions of these directions
follow the current nondominated distribution. Uniform vectors are used when the
adaptive construction is not activated or its fallback condition is met.

Each objective vector is associated with the direction of largest cosine
similarity, $a(i)=\arg\max_j
\vf_i^{\mathsf T}\vv_j/(\|\vf_i\|_2\|\vv_j\|_2)$, and the conventional
PBI value~\cite{ref:moead} along that direction is
\begin{equation}
\begin{aligned}
d_{1,i}&=(\vf_i-\vz^{*})^{\mathsf T}\vv_{a(i)}\big/\|\vv_{a(i)}\|_2,\\
d_{2,i}&=\bigl\|(\vf_i-\vz^{*})
        -d_{1,i}\,\vv_{a(i)}/\|\vv_{a(i)}\|_2\bigr\|_2,\\
g_i^{\mathrm{PBI}}&=d_{1,i}+\theta\,d_{2,i},
\end{aligned}
\label{eq:pbi}
\end{equation}
with penalty $\theta$. The continuous quality score is the decreasing transform
\begin{equation}
S_i=\frac{1}{1+g_i^{\mathrm{PBI}}}\in(0,1],
\label{eq:score-cont}
\end{equation}
so a larger $S_i$ indicates a smaller PBI value along the associated direction.
Unlike a class label, $S_i$ can distinguish solutions assigned the same binary
label.
Equations \eqref{eq:pbi}--\eqref{eq:score-cont} also fix the coordinate
convention of the implementation: the association is computed from the original
objective vectors, whereas the two distances are computed after translation by
$\vz^{*}$. All experiments in this paper use this single convention.

\subsubsection{Representative-Based Binary Group Label}
\label{sec:method:hpc:coarse}

The continuous score supplies a fine-grained ordering, while the binary label
adds a representative-relative quality boundary. This boundary is derived from a small set
$\mathcal{R}\subset\mathcal{P}$ of $k$ representative solutions, selected from
the current population by the radial-grid representative selection of
RSEA~\cite{ref:rsea}, which is also the rule used by REMO~\cite{ref:remo}. The
representatives are evaluated solutions that act as dynamic, population-derived
reference points. Their number increases with the objective dimension to provide
finer regional coverage as the many-objective geometry becomes richer.

Given $\mathcal{R}$, the binary group label applies the PBI partition of
REMO~\cite{ref:remo} with the representatives as reference points. Every
solution is associated with the representative $\vr_{b(i)}$ of largest cosine
similarity, the associated direction is
$\vw_{b(i)}=(\vr_{b(i)}-\vz^{*})/\|\vr_{b(i)}-\vz^{*}\|_2$, the two distances
$\hat d_{1,i},\hat d_{2,i}$ are formed as in \eqref{eq:pbi} with
$\vv_{a(i)}$ replaced by $\vw_{b(i)}$, and the binary label is
\begin{equation}
L_i=\mathbb{I}\!\left[
\frac{\hat d_{1,i}+\delta\,\hat d_{2,i}}
     {\|\vr_{b(i)}-\vz^{*}\|_2}\leq1\right]\in\{0,1\}.
\label{eq:label-bin}
\end{equation}
Normalising by $\|\vr_{b(i)}-\vz^{*}\|_2$ makes the associated representative
itself the threshold, so $L_i$ records on which side of the surface through that
representative a solution lies. Following REMO, the class-balance variable
$\delta$ is adjusted by bounded bisection so that the boundary maintains useful
occupancy on both sides across changing population states.

\subsubsection{Hybrid Score and Positive-Group Construction}
\label{sec:method:hpc:hybrid}

The two signals are combined by a weight that follows the consumption of the
expensive budget. With $t=FE/FE_{\max}$ and $\alpha_t=1-t$, the hybrid quality
score of solution $i$ is
\begin{equation}
H_i=\alpha_t S_i+(1-\alpha_t)L_i ,
\qquad
\alpha_t=1-t ,
\label{eq:fuse}
\end{equation}
and the positive group collects the highest-ranked fraction $r_g$ of the
population,
\begin{equation}
\mathcal{C}_1=\bigl\{\vx_i\;\big|\;
\operatorname{rank}_{\downarrow}(H_i)\leq\lceil r_gN\rceil\bigr\},
\qquad
\mathcal{C}_2=\mathcal{P}\setminus\mathcal{C}_1 .
\label{eq:group}
\end{equation}
With $r_g<1/2$ the non-positive group $\mathcal{C}_2$ provides a broad contrast
set, while the compact positive group concentrates pairwise supervision on the
highest-priority portion of the population.

The hybrid score connects quality assessment to relation learning through the
chain
\begin{equation*}
\begin{aligned}
&\text{quality score}\;\rightarrow\;\text{positive-group construction}\\
&\qquad\rightarrow\;\text{relation-pair construction},
\end{aligned}
\end{equation*}
so a change in the ranking matters insofar as it changes
which solutions enter $\mathcal{C}_1$ and therefore which ordered pairs the
relation model is trained on. The budget-dependent weighting connects
fine-grained quality ranking with representative-based group priority. As the
evaluation budget is consumed, the binary label receives increasing weight,
while $S_i$ refines the within-label order at each training step with $FE<FE_{\max}$.
Proposition~\ref{prop:order} characterises the resulting ordering.

\begin{proposition}[Ordering induced by the hybrid score]
\label{prop:order}
Let $S_i\in(0,1]$ as guaranteed by \eqref{eq:score-cont} whenever
$\vv_j\succeq\bm{0}$ and $\vf_i\succeq\vz^{*}$, so that
$g_i^{\mathrm{PBI}}\geq0$. Then
\begin{enumerate}[label=(\roman*),leftmargin=*,itemsep=1pt]
\item for any pair with $L_i=L_j$,
      $\;H_i-H_j=\alpha_t\,(S_i-S_j)$;
\item for any pair with $L_i=1$ and $L_j=0$,
      $\;\alpha_t\leq1/2\;\Rightarrow\;H_i>H_j$.
\end{enumerate}
\end{proposition}

\begin{proof}
(i) is immediate from \eqref{eq:fuse} because the term $(1-\alpha_t)L_i$
cancels. For (ii), if $\alpha_t=0$, then $H_i=1>0=H_j$.
For $0<\alpha_t\leq1/2$, $H_i=(1-\alpha_t)+\alpha_tS_i>1-\alpha_t$ because $S_i>0$,
while $H_j=\alpha_tS_j\leq\alpha_t$ because $S_j\leq1$; for
$\alpha_t\leq1/2$ we have $1-\alpha_t\geq\alpha_t$, hence $H_i>H_j$.
\end{proof}

Part~(i) explains why the continuous score is included: for $\alpha_t>0$,
solutions assigned the same binary label can still be distinguished when the
positive group is formed. Part~(ii) gives the schedule a concrete interpretation.
For $t<1/2$, the continuous score can affect ordering across the binary boundary;
for $1/2\leq t<1$, binary-positive solutions take priority and the continuous
score refines the order within each label. This transition combines fine-grained
ranking with a clear group priority as the budget is consumed.
Section~\ref{sec:exp:hpc} examines the resulting group composition and its
association with subsequent solution retention.

\subsection{Relation Learning on the Hybrid Groups}
\label{sec:method:relation}

HPDC-MaOEA retains the relation model of REMO~\cite{ref:remo} and supplies it
with the redesigned hybrid groups. Every ordered pair
of distinct solutions becomes one training sample whose input is the
concatenation $[\vx_i,\vx_j]\in\mathbb{R}^{2D}$ and whose label
\begin{equation}
y_{ij}=
\begin{cases}
+1, & \vx_i\in\mathcal{C}_1,\ \vx_j\in\mathcal{C}_2,\\
-1, & \vx_i\in\mathcal{C}_2,\ \vx_j\in\mathcal{C}_1,\\
0,  & \text{otherwise},
\end{cases}
\label{eq:pairlabel}
\end{equation}
with $\mathcal{C}_1,\mathcal{C}_2$ from \eqref{eq:group}, records the ordered
group relation rather than a Pareto comparison. Because $r_g<1/2$ makes the
same-group families larger than the cross-group ones, they are randomly
subsampled towards the cross-group count, so the three labels of
\eqref{eq:pairlabel} are balanced by construction rather than by a reweighted
loss.

We employ the same relation-learning architecture as REMO: min--max normalised
pair inputs, a class-balanced training/held-out split, and a feed-forward pattern
network ending in a softmax over the three labels.
Its output for a pair is written
$\bm{\pi}(\vx_a,\vx_b)=[\pi_{+1},\pi_0,\pi_{-1}]$, the predicted probabilities
that the first solution belongs to the higher group, that both belong to the
same group, and that the second belongs to the higher group. The held-out
misclassification rate $e_r$ is reused by the reliability-aware acquisition
rule in Section~\ref{sec:method:candidate:exp}.

For an unevaluated candidate $\vx$ the trained network is queried on the four
families of ordered comparisons $(\mathcal{C}_1,\vx)$, $(\vx,\mathcal{C}_1)$,
$(\mathcal{C}_2,\vx)$ and $(\vx,\mathcal{C}_2)$, that is on $2N$ pairs per
candidate. With $\mathbf{a},\mathbf{b},\mathbf{c},\mathbf{d}$ the family-wise
mean probability vectors in that order, the relation score aggregates the
evidence that $\vx$ is not inferior to the positive group and is superior to the
non-positive group, minus the opposite evidence,
\begin{equation}
R(\vx)=2\bigl[c_{-1}(\vx)+d_{+1}(\vx)-a_{+1}(\vx)-b_{-1}(\vx)\bigr]
\in[-4,4],
\label{eq:relscore}
\end{equation}
the bound following from each probability vector summing to one. $R(\vx)$ is the
aggregate group-relative preference used to order unevaluated candidates.

\subsection{Dual-Mode Candidate Selection}
\label{sec:method:candidate}

The relation surrogate can rank thousands of predicted-promising candidates per
iteration, but spending the expensive budget on the top of that single ranking
concentrates the batch on whatever preference the current surrogate happens to
hold. In particular, the number of candidates passing a fixed cutoff depends on
the distribution of $R(\vx)$, which changes with the problem, search stage and
training set. The resulting batch size, and hence the number of evaluations
spent in an iteration, can vary even though the cutoff is unchanged
(Section~\ref{sec:exp:candidate}). Moreover, $R(\vx)$ expresses only a
candidate's predicted position relative to the two training groups.

HPDC-MaOEA responds with two design decisions. First, ranking and expenditure
control are separated: quantile rules define thresholds relative to the current
score distribution, while an explicit cap $n_{\max}$ bounds the evaluation
batch. Scores determine candidate priority within this bound; the retained-set
size and remaining budget determine the feasible batch size. Second, two alternative ranking
criteria are made available --- an exploration-oriented mode
(Section~\ref{sec:method:candidate:exp}) and an indicator-oriented mode
(Section~\ref{sec:method:candidate:ind}) --- and one of them is drawn per
iteration (Section~\ref{sec:method:candidate:mix}). The two modes provide two
alternative criteria for allocating expensive evaluations while keeping the
per-iteration expenditure explicit and bounded.

Both modes rank the same kind of candidate pool, produced by a relation-guided
inner search: genetic variation is applied to the union of the current
population and the representatives $\mathcal{R}$, the relation score
\eqref{eq:relscore} ranks the offspring, the best of them are retained as
parents of the next inner generation together with $\mathcal{R}$, and the loop
stops once the cumulative number of generated candidates reaches the limit
$g_{\max}$. All candidates generated along the way are accumulated and
deduplicated into the pool $\mathcal{A}$. Mode-specific relation aggregation
shapes the inner search and produces a criterion-specific pool, while the common
pool-construction rule and batch bound preserve a consistent expenditure
interface. Section~\ref{sec:exp:candidate} examines how the two criteria change
the selected batches and the resulting search performance.

\subsubsection{Exploration-Oriented Mode}
\label{sec:method:candidate:exp}

This mode supplements the relation score with two quantities the score alone
does not carry: how sharply the classifier separates the classes around a
candidate, and how far the candidate lies from the rest of the batch in
decision space. Let $\mathcal{O}(\vx)$ collect the $2N$ ordered comparisons of
$\vx$ used in \eqref{eq:relscore}. The prediction ambiguity and the
exploration-oriented acquisition score are
\begin{equation}
\begin{aligned}
U(\vx)&=1-\frac{1}{|\mathcal{O}(\vx)|}
\sum_{(\vx_a,\vx_b)\in\mathcal{O}(\vx)}\;
\max_{y}\;\pi_y(\vx_a,\vx_b),\\
A_{\mathrm{exp}}(\vx)&=\widetilde R(\vx)+\lambda_t\,\widetilde U(\vx),
\end{aligned}
\label{eq:expscore}
\end{equation}
where $\widetilde{(\cdot)}$ denotes min--max normalisation over the pool. A
large $U(\vx)$ means the network produces flatter class distributions around
that candidate. The resulting softmax-concentration signal supplies a relative
acquisition order over the current pool. In this mode the same sharpness
$\max_y\pi_y$ additionally weights the family-wise averages entering
\eqref{eq:relscore}, so that sharper pairs contribute more to the aggregated
score while preserving the semantics of the three relation classes.

The weight of the ambiguity term is reduced according to the search stage and
the error of the model whose ambiguity is being sampled,
\begin{equation}
\lambda_t=\lambda_0(1-t)\max\!\left(0,\;1-\frac{e_r}{e_{\max}}\right),
\label{eq:lambda}
\end{equation}
with $e_r$ the held-out pairwise error of
Section~\ref{sec:method:relation} and $e_{\max}$ a fixed threshold. The
reliability factor automatically suppresses ambiguity guidance as the
held-out error approaches the threshold, while the progress factor shifts the
acquisition towards relation quality as the budget is consumed.

Candidates at or above the $q_{\mathrm{keep}}$ quantile of $A_{\mathrm{exp}}$
form the retained set
\begin{equation}
\mathcal{H}^{\mathrm{exp}}=\Bigl\{\vx\in\mathcal{A}\;\Bigm|\;
A_{\mathrm{exp}}(\vx)\geq
Q_{q_{\mathrm{keep}}}\bigl(A_{\mathrm{exp}}(\mathcal{A})\bigr)\Bigr\},
\label{eq:prefilter-exp}
\end{equation}
completed with the highest-scoring candidates if fewer than $n_{\min}$ qualify.
The batch is then assembled greedily from $\mathcal{H}^{\mathrm{exp}}$: the
first member maximises $A_{\mathrm{exp}}$, and each subsequent member maximises
\begin{equation}
\begin{split}
A_{\mathrm{batch}}(\vx\mid\mathcal{S})
&=w\,\widehat A_{\mathrm{exp}}(\vx)
+(1-w)\,\widehat d(\vx,\mathcal{S}),\\
d(\vx,\mathcal{S})&=\min_{\vz\in\mathcal{S}}\|\vx-\vz\|_2,
\end{split}
\label{eq:greedy}
\end{equation}
where $\widehat{(\cdot)}$ is renormalised over the remaining candidates at each
step and $w>1/2$. Quality remains the primary criterion, while decision-space
distance distributes the expensive batch across distinct candidate regions.

\subsubsection{Indicator-Based Reranking Mode}
\label{sec:method:candidate:ind}

The second mode retains the relation model for coarse candidate filtering and
introduces a different quality criterion only for final reranking. Following
PIEA~\cite{ref:piea}, an SDE-based fitness~\cite{ref:sde} is computed for each
evaluated solution using the generalised $L_p$ front-shape estimate and its
finite convergence fallback. Because this fitness is available only for
evaluated solutions, an RBF-SVR model approximates it in the decision space,
\begin{equation}
\widehat I(\vx)=\mathcal{M}_I(\vx;\mathcal{D}),
\qquad
\mathcal{D}=\{(\vx_i,I_i)\}_{i=1}^{N},
\label{eq:svr}
\end{equation}
and reranks the pool. Selection then follows two stages. The relation score
\eqref{eq:relscore} first retains the top fraction $q_{\mathrm{rel}}$ of
$\mathcal{A}$, after which $\widehat I$ reranks this reduced set. Candidates at
or above its $q_{\mathrm{ind}}$ quantile form $\mathcal{H}^{\mathrm{ind}}$ and
are completed with the highest-ranked candidates if fewer than $n_{\min}$
qualify. The final batch contains the best $n_b$ members under $\widehat I$.
Thus, the contribution of this mode is the two-stage use of relation-based
filtering followed by indicator-based reranking.

\subsubsection{Probabilistic Mode Switching}
\label{sec:method:candidate:mix}

Let $m\in\{\mathrm{exp},\mathrm{ind}\}$ denote the mode of the current
iteration. With $u\sim U[0,1)$,
\begin{equation}
m=
\begin{cases}
\mathrm{ind}, & \text{if the indicator surrogate is available}\\
              & \text{and } u<p_{\mathrm{mix}},\\[2pt]
\mathrm{exp}, & \text{otherwise},
\end{cases}
\label{eq:mode}
\end{equation}
with $p_{\mathrm{mix}}=0.5$. Stochastic switching assigns one intact criterion
to each iteration, preserving the distinct selection semantics of the two modes
and giving them equal selection probability when the indicator surrogate is
available, without requiring their scores to be combined into a single ranking.

Whichever mode is drawn, the number of expensive evaluations is
\begin{equation}
n_b=\min\bigl(n_{\max},\,|\mathcal{H}^{m}|\bigr),
\label{eq:batch}
\end{equation}
truncated further by the remaining budget $FE_{\max}-FE$. Before this cap is
applied, the retained set is completed to $n_{\min}$ whenever the pool permits.
Together, the quantile rules and \eqref{eq:batch} convert surrogate scores into
a stable rank-and-budget interface. A direct variation step supplies candidates
when the inner pool is empty. Algorithm~\ref{alg:candidate} summarises the
module.

\begin{algorithm}[t]
\caption{Dual-mode candidate selection}
\label{alg:candidate}
\begin{algorithmic}[1]
\Require population $\mathcal{P}$, representatives $\mathcal{R}$, relation and
         indicator surrogates, mode $m$, progress $t$, held-out error $e_r$
\Ensure candidate batch $\mathcal{S}$
\State $\mathcal{A}\gets\emptyset$;
       $\mathcal{Q}\gets$ offspring of $\mathcal{P}\cup\mathcal{R}$
\While{$|\mathcal{A}|<g_{\max}$ and $\mathcal{Q}\neq\emptyset$}
  \State $\mathcal{A}\gets\mathcal{A}\cup\mathcal{Q}$
  \State Score $\mathcal{Q}$ by the relation surrogate, \eqref{eq:relscore},
         and keep the best of them as parents
  \State $\mathcal{Q}\gets$ offspring of the parents and $\mathcal{R}$
\EndWhile
\State Remove duplicates from $\mathcal{A}$
\If{$m=\mathrm{ind}$}
  \State Retain the top fraction $q_{\mathrm{rel}}$ of $\mathcal{A}$ by
         \eqref{eq:relscore}
  \State Rerank the retained subset by $\widehat I$ of \eqref{eq:svr} and keep
         its top quantile as $\mathcal{H}^{\mathrm{ind}}$
  \State $\mathcal{S}\gets$ the $n_b$ best members of
         $\mathcal{H}^{\mathrm{ind}}$ under $\widehat I$
\Else
  \State Compute $A_{\mathrm{exp}}$ by \eqref{eq:expscore}--\eqref{eq:lambda}
         and form $\mathcal{H}^{\mathrm{exp}}$ by \eqref{eq:prefilter-exp}
  \State $\mathcal{S}\gets$ greedy batch of $n_b$ members of
         $\mathcal{H}^{\mathrm{exp}}$ by \eqref{eq:greedy}
\EndIf
\State Bound $n_b$ by \eqref{eq:batch} and by the remaining budget
\State \Return $\mathcal{S}$
\end{algorithmic}
\end{algorithm}

\subsection{Computational Complexity}
\label{sec:method:complexity}

Let $N$ be the population size, $N_v=O(N)$ the number of directions, $N_c$ the
size of the deduplicated candidate pool, $D$ the number of decision variables
and $M$ the number of objectives.

In the grouping module, association and the continuous PBI computation
\eqref{eq:pbi}--\eqref{eq:score-cont} cost $O(NN_vM)$, the direction set
\eqref{eq:direction} costs one normalisation of the nondominated vectors, the
bounded bisection of \eqref{eq:label-bin} costs $O(B_\delta NkM)$ with
$B_\delta$ fixed by the bracket and the tolerance, and representative selection
adds nondominated sorting and radial-projection distances. With $N_v=O(N)$ and
$k=O(N)$ the module is dominated by $O(N^2M)$, and generating the relation pairs
adds $O(N^2D)$.

In the candidate module, scoring one candidate requires $2N$ network inputs of
dimension $2D$, so the pool costs $O(N_cND)$ in data construction and
prediction, which dominates the module. The indicator fitness costs $O(N^2M)$,
the indicator-oriented mode adds one SVR fit on the $N$ evaluated solutions plus
predictions on the screened subset, and the greedy construction \eqref{eq:greedy}
costs $O(n_{\max}N_cD)$. The per-iteration cost is therefore $O(N_cND+N^2M)$.
Neither module consumes expensive evaluations, so the total number of real
evaluations remains $FE_{\max}$ regardless of the mode sequence; the wall-clock
share of the two modules is reported in Section~\ref{sec:exp:sensitivity}.


\section{Experimental Studies}
\label{sec:exp}

The experiments evaluated overall optimisation performance, the contribution
of the two modules, and the group and batch properties associated with their
design. The main comparisons used 10 and 20 objectives. Additional objective
counts and separate mechanism studies provided complementary evidence.

\subsection{Experimental Settings}
\label{sec:exp:settings}

\paragraph{Benchmarks and evaluation budgets.}
The main test suite comprised DTLZ1--DTLZ7 and WFG1--WFG9. Both the 10- and
20-objective comparisons used a population setting of $N=100$ and a budget of
$FE_{\max}=300$ real evaluations. The recorded decision dimension was $D=30$,
except for WFG3, which used $D=31$. Additional results covered $M=3,5,8$ on the
same 16 problems and $M=12$ on WFG1--WFG9. Table~\ref{tab:exp:settings}
distinguishes these settings from the mechanism experiments. In particular,
the grouping studies used 500 evaluations, whereas the candidate-value probe
used 300. Their results were analysed within their respective protocols.

\begin{table*}[t]
\centering
\caption{Experimental settings of the retained result sets. A dash denotes unverified or unavailable metadata.}
\label{tab:exp:settings}
\footnotesize
\setlength{\tabcolsep}{3pt}
\renewcommand{\arraystretch}{1.12}
\begin{tabular}{lccccc}
\toprule
Result set & $M$ & $D$ & $N$ & $FE_{\max}$ & Runs/config. \\
\midrule
Main comparison & 10, 20 & 30 (31) & 100 & 300 & -- \\
Additional full suite & 3, 5 & 10 & -- & -- & -- \\
Additional full suite & 8 & 10 (11) & 100 & 300 & -- \\
WFG-only comparison & 12 & 30 (31) & 100 & 300 & -- \\
Two-module ablation & 10, 20 & 30 (31) & 100 & 300 & -- \\
Snapshot/group audit & 10, 20 & 30 (31) & 100 & 500 & 5/behaviour \\
Good-group precision & 10, 20 & 30 (31) & 100 & 500 & 25 \\
Candidate-value probe & 20 & 30 (31) & 100 & 300 & 10/policy \\
\bottomrule
\end{tabular}
\par\smallskip
\begin{minipage}{\linewidth}
\scriptsize Parenthesised decision dimensions apply to WFG3. The snapshot audit has two behaviours per configuration; the candidate probe has five policies. The good-group precision results cover the original five-problem study. Spreadsheet run counts are intentionally left unfilled.
\end{minipage}
\end{table*}

\paragraph{Algorithms and indicators.}
HPDC-MaOEA was compared with REMO, PIEA, MCEAD, R2AEA, KRVEA and CSEA at both
$M=10$ and $M=20$, with PCSAEA additionally available at $M=10$ and HES-EA at
$M=20$. REMO, PIEA, MCEAD and KRVEA formed the common subset across the
additional objective counts. All available baselines were retained in each
main comparison. The implementation recorded as
\texttt{REMO\_new2\_AdaMaO\_SDEOnly\_UniformMix\_Original} is denoted
HPDC-MaOEA in the performance tables.
The primary indicator was inverted generational distance (IGD), for which
smaller values indicate better approximation quality. Entries report the
recorded mean and standard deviation; the lowest mean in each row is bold.
The separate routing study also supplied IGD$^+$ results at $M=20$.

\paragraph{Statistical reporting.}
The main comparison exports identify their pairwise annotations as Wilcoxon
rank-sum results. We preserve their direction: $+$ denotes a comparator
reported as better than HPDC-MaOEA, $-$ a comparator reported as worse, and
$=$ no detected difference. Summary counts are expressed from HPDC-MaOEA's
perspective as wins/ties/losses (W/T/L), so an exported $-$ contributes a win.
Counts were reconstructed from the individual cells. These annotations are
reported as supplied; no new significance test was inferred from means and
standard deviations. In the mechanism studies, repeated checkpoints were
first aggregated within each run. The grouping complementarity tests used
25 runs per problem--objective--stage cell and Holm adjustment over 40 cells.
The candidate probe is reported through run-level descriptive statistics.

\TODO{Complete the independent-run counts, original test settings and
multiple-comparison policy for the spreadsheet experiments; recover the
missing $N$ and FE settings for $M=3,5$, and the execution platform and IGD
reference-set specifications. Verify the identical PCSAEA/KRVEA entries for
DTLZ1 at $M=10$ before finalising the corresponding comparisons.}

\subsection{Comparison with State-of-the-Art Algorithms}
\label{sec:exp:comparison}

HPDC-MaOEA obtained more reported wins than losses across the available
comparators in both main settings. The aggregate W/T/L counts were
73/18/21 at $M=10$ and 70/21/21 at $M=20$, each over 112 problem--comparator
pairs (Tables~\ref{tab:exp:main10} and~\ref{tab:exp:main20}). Against REMO,
the counts were 9/5/2 and 8/6/2, respectively; against CSEA, they were
12/3/1 and 11/4/1. The comparison with PIEA was more problem dependent,
with 6/3/7 at $M=10$ and 9/2/5 at $M=20$.

The strongest common improvements over REMO occurred on DTLZ7 and several
WFG problems. On DTLZ7, mean IGD decreased from 21.252 to 12.122 at $M=10$
and from 38.239 to 18.154 at $M=20$, reductions of 43.0\% and 52.5\%.
HPDC-MaOEA also attained the lowest mean among the available algorithms on
WFG7 in both settings, and on WFG1 and WFG2 at $M=10$. These results show
competitive performance across several many-objective instances under a
300-evaluation budget.

The problemwise results also distinguish this advantage from uniform
dominance. All seven comparators had lower reported IGD on WFG3 in both
main settings. R2AEA achieved lower mean IGD on DTLZ7 at both objective
counts, while PIEA and MCEAD were stronger on DTLZ5 and DTLZ6. Across the
four common baselines, the median relative IGD reduction was 10.49\% at
$M=10$ and 9.73\% at $M=20$; the corresponding arithmetic means were
3.37\% and $-0.17\%$. Relative reduction was calculated per
problem--comparator pair as $(\mathrm{IGD}_{b}-\mathrm{IGD}_{H})/
\mathrm{IGD}_{b}$, using the exported means. The difference between the
median and mean reflects the uneven magnitude of the problemwise effects.

\begin{table*}[t]
\centering
\caption{IGD comparison on 16 10-objective problems with 300 evaluations.}
\label{tab:exp:main10}
\scriptsize
\setlength{\tabcolsep}{3pt}
\renewcommand{\arraystretch}{1.12}
\begin{tabular}{lcccccccc}
\toprule
Problem & REMO & PIEA & MCEAD & R2AEA & PCSAEA & KRVEA & CSEA & HPDC-MaOEA \\
\midrule
DTLZ1 & \shortstack{$\mathbf{2.0748e+2}^{=}$\\$(4.35e+1)$} & \shortstack{$2.4207e+2^{=}$\\$(7.32e+1)$} & \shortstack{$2.3531e+2^{=}$\\$(4.54e+1)$} & \shortstack{$2.7206e+2^{-}$\\$(2.49e+1)$} & \shortstack{$3.6855e+2^{-\dagger}$\\$(2.80e+1)$} & \shortstack{$3.6855e+2^{-\dagger}$\\$(2.80e+1)$} & \shortstack{$2.2281e+2^{=}$\\$(4.96e+1)$} & \shortstack{$2.2944e+2$\\$(3.62e+1)$} \\
DTLZ2 & \shortstack{$1.0902e+0^{=}$\\$(9.26e-2)$} & \shortstack{$\mathbf{9.6871e-1}^{+}$\\$(5.69e-2)$} & \shortstack{$1.0974e+0^{=}$\\$(6.08e-2)$} & \shortstack{$1.4611e+0^{-}$\\$(5.68e-2)$} & \shortstack{$1.5378e+0^{-}$\\$(6.87e-2)$} & \shortstack{$1.5416e+0^{-}$\\$(6.16e-2)$} & \shortstack{$1.3229e+0^{-}$\\$(9.81e-2)$} & \shortstack{$1.0946e+0$\\$(8.88e-2)$} \\
DTLZ3 & \shortstack{$8.0813e+2^{+}$\\$(1.50e+2)$} & \shortstack{$\mathbf{5.5459e+2}^{+}$\\$(1.06e+2)$} & \shortstack{$5.7499e+2^{+}$\\$(1.31e+2)$} & \shortstack{$5.6353e+2^{+}$\\$(6.65e+1)$} & \shortstack{$1.3393e+3^{-}$\\$(8.67e+1)$} & \shortstack{$1.3574e+3^{-}$\\$(1.10e+2)$} & \shortstack{$9.0082e+2^{=}$\\$(1.72e+2)$} & \shortstack{$9.3919e+2$\\$(1.40e+2)$} \\
DTLZ4 & \shortstack{$\mathbf{1.1487e+0}^{=}$\\$(1.05e-1)$} & \shortstack{$1.2452e+0^{=}$\\$(4.05e-2)$} & \shortstack{$1.2487e+0^{-}$\\$(1.91e-2)$} & \shortstack{$1.5944e+0^{-}$\\$(7.20e-2)$} & \shortstack{$1.7339e+0^{-}$\\$(5.89e-2)$} & \shortstack{$1.7361e+0^{-}$\\$(6.76e-2)$} & \shortstack{$1.2198e+0^{=}$\\$(8.16e-2)$} & \shortstack{$1.1814e+0$\\$(1.12e-1)$} \\
DTLZ5 & \shortstack{$6.2052e-1^{=}$\\$(1.47e-1)$} & \shortstack{$\mathbf{2.9016e-1}^{+}$\\$(3.54e-2)$} & \shortstack{$3.7432e-1^{+}$\\$(6.73e-2)$} & \shortstack{$9.8650e-1^{-}$\\$(8.39e-2)$} & \shortstack{$9.8455e-1^{-}$\\$(6.03e-2)$} & \shortstack{$9.9808e-1^{-}$\\$(4.84e-2)$} & \shortstack{$7.6415e-1^{-}$\\$(1.07e-1)$} & \shortstack{$6.6049e-1$\\$(1.47e-1)$} \\
DTLZ6 & \shortstack{$1.6378e+1^{-}$\\$(7.30e-1)$} & \shortstack{$\mathbf{5.4851e+0}^{+}$\\$(1.81e+0)$} & \shortstack{$8.9744e+0^{+}$\\$(1.77e+0)$} & \shortstack{$1.3116e+1^{+}$\\$(1.22e+0)$} & \shortstack{$1.8023e+1^{-}$\\$(1.70e-1)$} & \shortstack{$1.8042e+1^{-}$\\$(1.58e-1)$} & \shortstack{$1.6064e+1^{-}$\\$(6.81e-1)$} & \shortstack{$1.4463e+1$\\$(1.10e+0)$} \\
DTLZ7 & \shortstack{$2.1252e+1^{-}$\\$(2.44e+0)$} & \shortstack{$1.8588e+1^{-}$\\$(5.79e+0)$} & \shortstack{$2.0445e+1^{-}$\\$(4.67e+0)$} & \shortstack{$\mathbf{6.2038e+0}^{+}$\\$(1.22e+0)$} & \shortstack{$2.9100e+1^{-}$\\$(1.69e+0)$} & \shortstack{$2.8736e+1^{-}$\\$(1.88e+0)$} & \shortstack{$2.7611e+1^{-}$\\$(2.24e+0)$} & \shortstack{$1.2122e+1$\\$(2.80e+0)$} \\
WFG1 & \shortstack{$3.2354e+0^{=}$\\$(8.74e-2)$} & \shortstack{$3.3433e+0^{-}$\\$(4.96e-2)$} & \shortstack{$3.4587e+0^{-}$\\$(3.65e-2)$} & \shortstack{$3.1986e+0^{=}$\\$(3.20e-2)$} & \shortstack{$3.4196e+0^{-}$\\$(1.68e-2)$} & \shortstack{$3.4256e+0^{-}$\\$(2.10e-2)$} & \shortstack{$3.2723e+0^{-}$\\$(6.63e-2)$} & \shortstack{$\mathbf{3.1933e+0}$\\$(2.79e-2)$} \\
WFG2 & \shortstack{$3.5108e+0^{-}$\\$(1.08e+0)$} & \shortstack{$2.3710e+0^{-}$\\$(4.69e-1)$} & \shortstack{$3.4836e+0^{-}$\\$(9.77e-1)$} & \shortstack{$2.2469e+0^{-}$\\$(1.13e-1)$} & \shortstack{$3.5949e+0^{-}$\\$(6.30e-1)$} & \shortstack{$3.8780e+0^{-}$\\$(6.00e-1)$} & \shortstack{$3.5554e+0^{-}$\\$(1.03e+0)$} & \shortstack{$\mathbf{2.1211e+0}$\\$(6.82e-1)$} \\
WFG3 & \shortstack{$1.5188e+0^{+}$\\$(1.01e-1)$} & \shortstack{$\mathbf{1.1823e+0}^{+}$\\$(1.49e-1)$} & \shortstack{$1.2551e+0^{+}$\\$(4.56e-2)$} & \shortstack{$1.5946e+0^{+}$\\$(4.99e-2)$} & \shortstack{$1.6472e+0^{+}$\\$(3.93e-2)$} & \shortstack{$1.6496e+0^{+}$\\$(4.64e-2)$} & \shortstack{$1.5723e+0^{+}$\\$(8.05e-2)$} & \shortstack{$1.7006e+0$\\$(5.87e-2)$} \\
WFG4 & \shortstack{$8.1770e+0^{-}$\\$(9.93e-1)$} & \shortstack{$\mathbf{5.7766e+0}^{+}$\\$(5.86e-1)$} & \shortstack{$7.1626e+0^{=}$\\$(8.70e-1)$} & \shortstack{$5.7901e+0^{+}$\\$(3.16e-1)$} & \shortstack{$9.1390e+0^{-}$\\$(4.55e-1)$} & \shortstack{$9.2700e+0^{-}$\\$(3.98e-1)$} & \shortstack{$9.0803e+0^{-}$\\$(8.62e-1)$} & \shortstack{$6.6311e+0$\\$(6.38e-1)$} \\
WFG5 & \shortstack{$6.9444e+0^{-}$\\$(5.14e-1)$} & \shortstack{$\mathbf{5.7544e+0}^{=}$\\$(2.59e-1)$} & \shortstack{$6.6193e+0^{-}$\\$(3.95e-1)$} & \shortstack{$6.5834e+0^{-}$\\$(4.62e-1)$} & \shortstack{$7.0750e+0^{-}$\\$(1.97e-1)$} & \shortstack{$6.9475e+0^{-}$\\$(2.47e-1)$} & \shortstack{$7.1022e+0^{-}$\\$(4.09e-1)$} & \shortstack{$5.9732e+0$\\$(4.00e-1)$} \\
WFG6 & \shortstack{$6.4255e+0^{-}$\\$(4.79e-1)$} & \shortstack{$5.9360e+0^{-}$\\$(4.03e-1)$} & \shortstack{$7.2512e+0^{-}$\\$(6.73e-1)$} & \shortstack{$\mathbf{5.2560e+0}^{=}$\\$(1.32e-1)$} & \shortstack{$7.3710e+0^{-}$\\$(3.74e-1)$} & \shortstack{$7.3224e+0^{-}$\\$(3.02e-1)$} & \shortstack{$7.4553e+0^{-}$\\$(5.14e-1)$} & \shortstack{$5.2773e+0$\\$(2.01e-1)$} \\
WFG7 & \shortstack{$6.2812e+0^{-}$\\$(7.42e-1)$} & \shortstack{$6.9139e+0^{-}$\\$(5.04e-1)$} & \shortstack{$8.1581e+0^{-}$\\$(6.29e-1)$} & \shortstack{$6.4875e+0^{-}$\\$(4.09e-1)$} & \shortstack{$7.8445e+0^{-}$\\$(3.81e-1)$} & \shortstack{$7.8642e+0^{-}$\\$(4.20e-1)$} & \shortstack{$7.3904e+0^{-}$\\$(6.15e-1)$} & \shortstack{$\mathbf{5.2969e+0}$\\$(2.57e-1)$} \\
WFG8 & \shortstack{$7.0559e+0^{-}$\\$(5.06e-1)$} & \shortstack{$6.3573e+0^{-}$\\$(4.99e-1)$} & \shortstack{$7.4671e+0^{-}$\\$(4.72e-1)$} & \shortstack{$\mathbf{5.5601e+0}^{=}$\\$(1.19e-1)$} & \shortstack{$7.5364e+0^{-}$\\$(2.73e-1)$} & \shortstack{$7.4984e+0^{-}$\\$(3.01e-1)$} & \shortstack{$7.7877e+0^{-}$\\$(3.46e-1)$} & \shortstack{$5.6930e+0$\\$(2.66e-1)$} \\
WFG9 & \shortstack{$7.6476e+0^{-}$\\$(7.61e-1)$} & \shortstack{$\mathbf{5.7930e+0}^{+}$\\$(3.72e-1)$} & \shortstack{$7.0825e+0^{-}$\\$(6.96e-1)$} & \shortstack{$6.4412e+0^{=}$\\$(3.15e-1)$} & \shortstack{$7.7853e+0^{-}$\\$(4.44e-1)$} & \shortstack{$7.7844e+0^{-}$\\$(3.65e-1)$} & \shortstack{$8.1835e+0^{-}$\\$(3.07e-1)$} & \shortstack{$6.3788e+0$\\$(6.11e-1)$} \\
\midrule
W/T/L & 9/5/2 & 6/3/7 & 9/3/4 & 7/4/5 & 15/0/1 & 15/0/1 & 12/3/1 & -- \\
\bottomrule
\end{tabular}
\par\smallskip
\begin{minipage}{\linewidth}
\scriptsize Each cell contains the exported mean IGD and standard deviation on separate lines. Symbols describe the comparator relative to HPDC-MaOEA; W/T/L describes HPDC-MaOEA relative to the comparator. Bold denotes the smallest mean. $\dagger$: the two identical DTLZ1 entries are preserved from the source and require provenance verification.
\end{minipage}
\end{table*}

\begin{table*}[t]
\centering
\caption{IGD comparison on 16 20-objective problems with 300 evaluations.}
\label{tab:exp:main20}
\scriptsize
\setlength{\tabcolsep}{3pt}
\renewcommand{\arraystretch}{1.12}
\begin{tabular}{lcccccccc}
\toprule
Problem & REMO & PIEA & MCEAD & R2AEA & KRVEA & CSEA & HES-EA & HPDC-MaOEA \\
\midrule
DTLZ1 & \shortstack{$\mathbf{6.7408e+1}^{+}$\\$(1.93e+1)$} & \shortstack{$9.3984e+1^{-}$\\$(1.82e+1)$} & \shortstack{$7.6364e+1^{=}$\\$(3.61e+1)$} & \shortstack{$1.0965e+2^{-}$\\$(1.62e+1)$} & \shortstack{$1.5756e+2^{-}$\\$(1.98e+1)$} & \shortstack{$7.9813e+1^{=}$\\$(1.99e+1)$} & \shortstack{$1.6595e+2^{-}$\\$(1.85e+1)$} & \shortstack{$7.9755e+1$\\$(1.09e+1)$} \\
DTLZ2 & \shortstack{$\mathbf{1.1173e+0}^{=}$\\$(5.45e-2)$} & \shortstack{$1.1636e+0^{-}$\\$(2.17e-2)$} & \shortstack{$1.2196e+0^{-}$\\$(2.96e-2)$} & \shortstack{$1.3157e+0^{-}$\\$(4.47e-2)$} & \shortstack{$1.3465e+0^{-}$\\$(2.76e-2)$} & \shortstack{$1.1817e+0^{-}$\\$(5.61e-2)$} & \shortstack{$1.3478e+0^{-}$\\$(3.70e-2)$} & \shortstack{$1.1350e+0$\\$(4.89e-2)$} \\
DTLZ3 & \shortstack{$2.6646e+2^{=}$\\$(7.13e+1)$} & \shortstack{$2.5997e+2^{=}$\\$(1.15e+1)$} & \shortstack{$\mathbf{1.2168e+2}^{+}$\\$(1.12e+2)$} & \shortstack{$3.0721e+2^{=}$\\$(5.93e+1)$} & \shortstack{$5.1789e+2^{-}$\\$(7.55e+1)$} & \shortstack{$2.5514e+2^{=}$\\$(6.82e+1)$} & \shortstack{$5.4369e+2^{-}$\\$(8.04e+1)$} & \shortstack{$2.8568e+2$\\$(6.84e+1)$} \\
DTLZ4 & \shortstack{$\mathbf{1.0368e+0}^{=}$\\$(4.26e-2)$} & \shortstack{$1.1269e+0^{-}$\\$(6.03e-2)$} & \shortstack{$1.2852e+0^{-}$\\$(3.78e-2)$} & \shortstack{$1.2275e+0^{-}$\\$(5.62e-2)$} & \shortstack{$1.2558e+0^{-}$\\$(3.34e-2)$} & \shortstack{$1.0610e+0^{=}$\\$(4.81e-2)$} & \shortstack{$1.2647e+0^{-}$\\$(2.73e-2)$} & \shortstack{$1.0578e+0$\\$(5.34e-2)$} \\
DTLZ5 & \shortstack{$2.7068e-1^{=}$\\$(5.25e-2)$} & \shortstack{$1.5237e-1^{+}$\\$(2.26e-2)$} & \shortstack{$\mathbf{1.5161e-1}^{+}$\\$(4.26e-2)$} & \shortstack{$4.4572e-1^{-}$\\$(7.26e-2)$} & \shortstack{$4.1362e-1^{-}$\\$(4.70e-2)$} & \shortstack{$2.9000e-1^{=}$\\$(5.36e-2)$} & \shortstack{$4.1158e-1^{-}$\\$(3.61e-2)$} & \shortstack{$3.0345e-1$\\$(6.65e-2)$} \\
DTLZ6 & \shortstack{$7.3417e+0^{=}$\\$(6.33e-1)$} & \shortstack{$\mathbf{2.1640e+0}^{+}$\\$(1.05e+0)$} & \shortstack{$4.1169e+0^{+}$\\$(1.03e+0)$} & \shortstack{$5.6087e+0^{=}$\\$(1.98e+0)$} & \shortstack{$9.2234e+0^{-}$\\$(1.53e-1)$} & \shortstack{$7.4321e+0^{-}$\\$(7.00e-1)$} & \shortstack{$9.1974e+0^{-}$\\$(1.22e-1)$} & \shortstack{$6.6546e+0$\\$(9.39e-1)$} \\
DTLZ7 & \shortstack{$3.8239e+1^{-}$\\$(7.37e+0)$} & \shortstack{$1.8728e+1^{-}$\\$(9.86e+0)$} & \shortstack{$2.8155e+1^{-}$\\$(1.25e+1)$} & \shortstack{$\mathbf{1.4039e+1}^{+}$\\$(2.89e+0)$} & \shortstack{$4.9744e+1^{-}$\\$(5.36e+0)$} & \shortstack{$4.8818e+1^{-}$\\$(9.05e+0)$} & \shortstack{$4.8867e+1^{-}$\\$(4.77e+0)$} & \shortstack{$1.8154e+1$\\$(2.61e+0)$} \\
WFG1 & \shortstack{$5.6202e+0^{=}$\\$(8.82e-2)$} & \shortstack{$5.7081e+0^{-}$\\$(5.05e-2)$} & \shortstack{$5.7957e+0^{-}$\\$(4.51e-2)$} & \shortstack{$\mathbf{5.6059e+0}^{=}$\\$(4.65e-2)$} & \shortstack{$5.7218e+0^{-}$\\$(3.08e-2)$} & \shortstack{$5.6736e+0^{-}$\\$(1.06e-1)$} & \shortstack{$5.7151e+0^{-}$\\$(3.91e-2)$} & \shortstack{$5.6075e+0$\\$(7.93e-2)$} \\
WFG2 & \shortstack{$7.9480e+0^{-}$\\$(1.54e+0)$} & \shortstack{$6.0506e+0^{-}$\\$(1.05e+0)$} & \shortstack{$7.6391e+0^{-}$\\$(2.05e+0)$} & \shortstack{$\mathbf{4.6261e+0}^{+}$\\$(1.67e-1)$} & \shortstack{$8.3973e+0^{-}$\\$(9.72e-1)$} & \shortstack{$8.8363e+0^{-}$\\$(1.30e+0)$} & \shortstack{$8.8430e+0^{-}$\\$(1.37e+0)$} & \shortstack{$5.4393e+0$\\$(1.31e+0)$} \\
WFG3 & \shortstack{$2.2519e+0^{+}$\\$(2.32e-1)$} & \shortstack{$\mathbf{1.7073e+0}^{+}$\\$(1.73e-1)$} & \shortstack{$1.9974e+0^{+}$\\$(2.03e-1)$} & \shortstack{$2.6029e+0^{+}$\\$(1.15e-1)$} & \shortstack{$2.4955e+0^{+}$\\$(8.00e-2)$} & \shortstack{$2.3113e+0^{+}$\\$(1.84e-1)$} & \shortstack{$2.4573e+0^{+}$\\$(6.25e-2)$} & \shortstack{$2.7306e+0$\\$(1.52e-1)$} \\
WFG4 & \shortstack{$2.9066e+1^{-}$\\$(1.21e+0)$} & \shortstack{$2.0639e+1^{+}$\\$(3.25e+0)$} & \shortstack{$2.0456e+1^{+}$\\$(3.65e+0)$} & \shortstack{$\mathbf{1.7947e+1}^{+}$\\$(1.25e+0)$} & \shortstack{$2.9490e+1^{-}$\\$(7.24e-1)$} & \shortstack{$2.9600e+1^{-}$\\$(1.48e+0)$} & \shortstack{$2.9432e+1^{-}$\\$(9.66e-1)$} & \shortstack{$2.3927e+1$\\$(1.07e+0)$} \\
WFG5 & \shortstack{$2.5608e+1^{-}$\\$(1.19e+0)$} & \shortstack{$\mathbf{2.1972e+1}^{+}$\\$(9.12e-1)$} & \shortstack{$2.2931e+1^{=}$\\$(6.52e-1)$} & \shortstack{$2.3519e+1^{=}$\\$(1.28e+0)$} & \shortstack{$2.4741e+1^{-}$\\$(6.96e-1)$} & \shortstack{$2.5912e+1^{-}$\\$(9.45e-1)$} & \shortstack{$2.4824e+1^{-}$\\$(5.85e-1)$} & \shortstack{$2.2849e+1$\\$(1.62e+0)$} \\
WFG6 & \shortstack{$2.4041e+1^{-}$\\$(1.66e+0)$} & \shortstack{$2.2696e+1^{-}$\\$(1.82e+0)$} & \shortstack{$2.5327e+1^{-}$\\$(1.94e+0)$} & \shortstack{$\mathbf{1.5679e+1}^{+}$\\$(5.02e-1)$} & \shortstack{$2.5293e+1^{-}$\\$(5.92e-1)$} & \shortstack{$2.6268e+1^{-}$\\$(1.20e+0)$} & \shortstack{$2.5393e+1^{-}$\\$(7.75e-1)$} & \shortstack{$1.9097e+1$\\$(6.18e-1)$} \\
WFG7 & \shortstack{$2.4322e+1^{-}$\\$(1.10e+0)$} & \shortstack{$2.4579e+1^{-}$\\$(8.74e-1)$} & \shortstack{$2.7355e+1^{-}$\\$(1.36e+0)$} & \shortstack{$2.0597e+1^{=}$\\$(2.63e+0)$} & \shortstack{$2.6465e+1^{-}$\\$(7.60e-1)$} & \shortstack{$2.6430e+1^{-}$\\$(1.40e+0)$} & \shortstack{$2.5957e+1^{-}$\\$(7.48e-1)$} & \shortstack{$\mathbf{2.0592e+1}$\\$(7.61e-1)$} \\
WFG8 & \shortstack{$2.4953e+1^{-}$\\$(1.02e+0)$} & \shortstack{$2.3237e+1^{-}$\\$(1.24e+0)$} & \shortstack{$2.6026e+1^{-}$\\$(1.81e+0)$} & \shortstack{$\mathbf{1.6282e+1}^{+}$\\$(3.72e-1)$} & \shortstack{$2.5519e+1^{-}$\\$(4.89e-1)$} & \shortstack{$2.7053e+1^{-}$\\$(7.42e-1)$} & \shortstack{$2.5791e+1^{-}$\\$(6.48e-1)$} & \shortstack{$2.0738e+1$\\$(8.26e-1)$} \\
WFG9 & \shortstack{$2.6155e+1^{-}$\\$(1.38e+0)$} & \shortstack{$2.2604e+1^{=}$\\$(1.32e+0)$} & \shortstack{$2.3089e+1^{=}$\\$(1.39e+0)$} & \shortstack{$\mathbf{2.2569e+1}^{=}$\\$(1.27e+0)$} & \shortstack{$2.5820e+1^{-}$\\$(7.55e-1)$} & \shortstack{$2.7847e+1^{-}$\\$(8.01e-1)$} & \shortstack{$2.6166e+1^{-}$\\$(8.22e-1)$} & \shortstack{$2.3053e+1$\\$(1.30e+0)$} \\
\midrule
W/T/L & 8/6/2 & 9/2/5 & 8/3/5 & 4/6/6 & 15/0/1 & 11/4/1 & 15/0/1 & -- \\
\bottomrule
\end{tabular}
\par\smallskip
\begin{minipage}{\linewidth}
\scriptsize Each cell contains the exported mean IGD and standard deviation on separate lines. Symbols describe the comparator relative to HPDC-MaOEA; W/T/L describes HPDC-MaOEA relative to the comparator. Bold denotes the smallest mean.
\end{minipage}
\end{table*}

\paragraph{Additional objective counts.}
Table~\ref{tab:exp:objectives} summarises the common-baseline comparisons.
On the full 16-problem suite, the reported win fraction increased from
34.4\% at $M=3$ to 62.5\% at $M=20$. This is a descriptive comparison
across the recorded settings: the decision dimension also changes between
the lower- and higher-objective experiments. The $M=12$ result covered
only nine WFG problems and is therefore shown separately. It yielded
27/3/6 against the common baselines and 43/3/8 against all six available
baselines.

\begin{table}[t]
\centering
\caption{Common-baseline comparisons across objective counts.}
\label{tab:exp:objectives}
\scriptsize
\setlength{\tabcolsep}{2.5pt}
\renewcommand{\arraystretch}{1.12}
\begin{tabular}{lcccc}
\toprule
$M$ & Problems & W/T/L & Win fraction & Median reduction \\
\midrule
3 & 16 & 22/18/24 & 34.4\% & 2.06\% \\
5 & 16 & 31/18/15 & 48.4\% & 6.76\% \\
8 & 16 & 38/12/14 & 59.4\% & 7.04\% \\
10 & 16 & 39/11/14 & 60.9\% & 10.49\% \\
20 & 16 & 40/11/13 & 62.5\% & 9.73\% \\
\midrule
12 & 9 & 27/3/6 & 75.0\% & 15.66\% \\
\bottomrule
\end{tabular}
\par\smallskip
\begin{minipage}{\linewidth}
\scriptsize REMO, PIEA, MCEAD and KRVEA are included in every row. The $M=12$ row covers WFG1--WFG9 only. Settings and verification notes follow Tables~\ref{tab:exp:settings}--\ref{tab:exp:main20}.
\end{minipage}
\end{table}

\subsection{Ablation Study}
\label{sec:exp:ablation}

The module ablation compared REMO, REMO with hybrid PBI grouping (HPC),
REMO with the candidate module, and the full method. Because the number
of representatives also affects the host algorithm, the available
configurations included REMO and REMO+HPC at both $k=6$ and the paper
setting $k=15$ for $M=10$ or $k=30$ for $M=20$. The candidate-only
configuration used $k=6$. Tables~\ref{tab:exp:abl10}
and~\ref{tab:exp:abl20} give all six identified configurations. The
20-objective table contains 15 valid problems; its WFG4 row is left blank
because the source ablation export labels that instance $M=10$.

At a matched representative count, adding candidate selection to
REMO+HPC produced reported W/T/L counts of 3/12/1 at $M=10$ and 3/11/1
at $M=20$. On DTLZ7, mean IGD decreased from 17.305 to 12.122 and from
25.637 to 18.154, respectively. WFG9 improved at both objective counts;
the remaining reported wins were WFG7 at $M=10$ and WFG6 at $M=20$.
Thus, the matched comparison identifies gains on specific problems while
most differences were not detected by the exported pairwise tests.

The full method attained the lowest mean IGD in 7 of the 16
10-objective instances and 5 of the 15 valid 20-objective instances among
the six configurations. Its comparison with enlarged-$k$ REMO yielded
10/5/1 and 10/4/1. These latter contrasts measure the combined change
from both modules at a fixed representative count. In contrast, the
candidate-only $k=6$ configuration differs from the full method in both
grouping and $k$, and does not isolate grouping removal at the paper
setting.

\begin{table*}[t]
\centering
\caption{Module ablation at $M=10$ and 300 evaluations.}
\label{tab:exp:abl10}
\scriptsize
\setlength{\tabcolsep}{4pt}
\renewcommand{\arraystretch}{1.12}
\begin{tabular}{lcccccc}
\toprule
Problem & \shortstack{REMO\\$k=6$} & \shortstack{REMO\\$k=15$} & \shortstack{REMO+HPC\\$k=6$} & \shortstack{REMO+HPC\\$k=15$} & \shortstack{REMO+candidate\\$k=6$} & \shortstack{HPDC-MaOEA\\$k=15$} \\
\midrule
DTLZ1 & \shortstack{$2.0748e+2^{=}$\\$(4.35e+1)$} & \shortstack{$2.5572e+2^{-}$\\$(3.72e+1)$} & \shortstack{$\mathbf{1.9210e+2}^{+}$\\$(4.29e+1)$} & \shortstack{$2.2265e+2^{=}$\\$(2.72e+1)$} & \shortstack{$2.1327e+2^{=}$\\$(4.33e+1)$} & \shortstack{$2.2944e+2$\\$(3.62e+1)$} \\
DTLZ2 & \shortstack{$1.0902e+0^{=}$\\$(9.26e-2)$} & \shortstack{$1.1321e+0^{=}$\\$(8.14e-2)$} & \shortstack{$1.1179e+0^{=}$\\$(8.60e-2)$} & \shortstack{$1.1325e+0^{=}$\\$(1.03e-1)$} & \shortstack{$\mathbf{1.0434e+0}^{=}$\\$(1.08e-1)$} & \shortstack{$1.0946e+0$\\$(8.88e-2)$} \\
DTLZ3 & \shortstack{$8.0813e+2^{+}$\\$(1.50e+2)$} & \shortstack{$9.8158e+2^{=}$\\$(1.46e+2)$} & \shortstack{$\mathbf{7.4833e+2}^{+}$\\$(1.57e+2)$} & \shortstack{$9.4062e+2^{=}$\\$(1.36e+2)$} & \shortstack{$8.3489e+2^{+}$\\$(1.38e+2)$} & \shortstack{$9.3919e+2$\\$(1.40e+2)$} \\
DTLZ4 & \shortstack{$\mathbf{1.1487e+0}^{=}$\\$(1.05e-1)$} & \shortstack{$1.2536e+0^{=}$\\$(1.07e-1)$} & \shortstack{$1.1609e+0^{=}$\\$(1.08e-1)$} & \shortstack{$1.1720e+0^{=}$\\$(8.09e-2)$} & \shortstack{$1.1829e+0^{=}$\\$(1.08e-1)$} & \shortstack{$1.1814e+0$\\$(1.12e-1)$} \\
DTLZ5 & \shortstack{$6.2052e-1^{=}$\\$(1.47e-1)$} & \shortstack{$6.5860e-1^{=}$\\$(1.07e-1)$} & \shortstack{$\mathbf{5.2409e-1}^{+}$\\$(1.16e-1)$} & \shortstack{$6.0132e-1^{=}$\\$(1.35e-1)$} & \shortstack{$5.8917e-1^{=}$\\$(1.12e-1)$} & \shortstack{$6.6049e-1$\\$(1.47e-1)$} \\
DTLZ6 & \shortstack{$1.6378e+1^{-}$\\$(7.30e-1)$} & \shortstack{$1.6101e+1^{-}$\\$(6.03e-1)$} & \shortstack{$1.5155e+1^{=}$\\$(6.93e-1)$} & \shortstack{$1.5000e+1^{=}$\\$(5.75e-1)$} & \shortstack{$1.6196e+1^{-}$\\$(5.75e-1)$} & \shortstack{$\mathbf{1.4463e+1}$\\$(1.10e+0)$} \\
DTLZ7 & \shortstack{$2.1252e+1^{-}$\\$(2.44e+0)$} & \shortstack{$2.1381e+1^{-}$\\$(2.48e+0)$} & \shortstack{$1.7953e+1^{-}$\\$(3.03e+0)$} & \shortstack{$1.7305e+1^{-}$\\$(1.70e+0)$} & \shortstack{$1.4917e+1^{-}$\\$(1.86e+0)$} & \shortstack{$\mathbf{1.2122e+1}$\\$(2.80e+0)$} \\
WFG1 & \shortstack{$3.2354e+0^{=}$\\$(8.74e-2)$} & \shortstack{$3.2346e+0^{-}$\\$(5.80e-2)$} & \shortstack{$3.1879e+0^{=}$\\$(5.68e-2)$} & \shortstack{$3.2000e+0^{=}$\\$(3.80e-2)$} & \shortstack{$\mathbf{3.1734e+0}^{=}$\\$(6.75e-2)$} & \shortstack{$3.1933e+0$\\$(2.79e-2)$} \\
WFG2 & \shortstack{$3.5108e+0^{-}$\\$(1.08e+0)$} & \shortstack{$2.7019e+0^{-}$\\$(8.69e-1)$} & \shortstack{$2.7960e+0^{-}$\\$(8.43e-1)$} & \shortstack{$2.1621e+0^{=}$\\$(6.46e-1)$} & \shortstack{$3.0646e+0^{-}$\\$(1.04e+0)$} & \shortstack{$\mathbf{2.1211e+0}$\\$(6.82e-1)$} \\
WFG3 & \shortstack{$\mathbf{1.5188e+0}^{+}$\\$(1.01e-1)$} & \shortstack{$1.5502e+0^{+}$\\$(8.35e-2)$} & \shortstack{$1.5760e+0^{+}$\\$(8.86e-2)$} & \shortstack{$1.6399e+0^{+}$\\$(5.18e-2)$} & \shortstack{$1.6296e+0^{+}$\\$(6.34e-2)$} & \shortstack{$1.7006e+0$\\$(5.87e-2)$} \\
WFG4 & \shortstack{$8.1770e+0^{-}$\\$(9.93e-1)$} & \shortstack{$7.2123e+0^{-}$\\$(7.13e-1)$} & \shortstack{$8.6573e+0^{-}$\\$(1.21e+0)$} & \shortstack{$7.1851e+0^{=}$\\$(9.13e-1)$} & \shortstack{$7.5869e+0^{-}$\\$(9.62e-1)$} & \shortstack{$\mathbf{6.6311e+0}$\\$(6.38e-1)$} \\
WFG5 & \shortstack{$6.9444e+0^{-}$\\$(5.14e-1)$} & \shortstack{$5.9358e+0^{=}$\\$(4.19e-1)$} & \shortstack{$6.6419e+0^{-}$\\$(6.91e-1)$} & \shortstack{$\mathbf{5.8962e+0}^{=}$\\$(3.37e-1)$} & \shortstack{$6.4368e+0^{-}$\\$(5.01e-1)$} & \shortstack{$5.9732e+0$\\$(4.00e-1)$} \\
WFG6 & \shortstack{$6.4255e+0^{-}$\\$(4.79e-1)$} & \shortstack{$5.4737e+0^{-}$\\$(2.95e-1)$} & \shortstack{$6.4196e+0^{-}$\\$(5.05e-1)$} & \shortstack{$5.3355e+0^{=}$\\$(3.10e-1)$} & \shortstack{$5.7033e+0^{-}$\\$(2.69e-1)$} & \shortstack{$\mathbf{5.2773e+0}$\\$(2.01e-1)$} \\
WFG7 & \shortstack{$6.2812e+0^{-}$\\$(7.42e-1)$} & \shortstack{$5.7119e+0^{-}$\\$(4.43e-1)$} & \shortstack{$6.6489e+0^{-}$\\$(6.28e-1)$} & \shortstack{$5.4193e+0^{-}$\\$(2.07e-1)$} & \shortstack{$5.7946e+0^{-}$\\$(3.58e-1)$} & \shortstack{$\mathbf{5.2969e+0}$\\$(2.57e-1)$} \\
WFG8 & \shortstack{$7.0559e+0^{-}$\\$(5.06e-1)$} & \shortstack{$6.0420e+0^{-}$\\$(3.51e-1)$} & \shortstack{$6.7588e+0^{-}$\\$(4.65e-1)$} & \shortstack{$\mathbf{5.6556e+0}^{=}$\\$(1.43e-1)$} & \shortstack{$6.3053e+0^{-}$\\$(3.64e-1)$} & \shortstack{$5.6930e+0$\\$(2.66e-1)$} \\
WFG9 & \shortstack{$7.6476e+0^{-}$\\$(7.61e-1)$} & \shortstack{$6.9496e+0^{-}$\\$(4.45e-1)$} & \shortstack{$7.6437e+0^{-}$\\$(6.71e-1)$} & \shortstack{$6.7894e+0^{-}$\\$(6.20e-1)$} & \shortstack{$7.1350e+0^{-}$\\$(5.95e-1)$} & \shortstack{$\mathbf{6.3788e+0}$\\$(6.11e-1)$} \\
\midrule
W/T/L & 9/5/2 & 10/5/1 & 8/4/4 & 3/12/1 & 9/5/2 & -- \\
Best mean & 2 & 0 & 3 & 2 & 2 & 7 \\
\bottomrule
\end{tabular}
\par\smallskip
\begin{minipage}{\linewidth}
\scriptsize Cells report mean IGD (standard deviation). Symbols and W/T/L follow the main tables. HPC denotes hybrid PBI grouping. The enlarged-$k$ HPC-only configuration isolates addition of candidate selection when compared with the full method.
\end{minipage}
\end{table*}

\begin{table*}[t]
\centering
\caption{Module ablation at $M=20$ and 300 evaluations.}
\label{tab:exp:abl20}
\scriptsize
\setlength{\tabcolsep}{4pt}
\renewcommand{\arraystretch}{1.12}
\begin{tabular}{lcccccc}
\toprule
Problem & \shortstack{REMO\\$k=6$} & \shortstack{REMO\\$k=30$} & \shortstack{REMO+HPC\\$k=6$} & \shortstack{REMO+HPC\\$k=30$} & \shortstack{REMO+candidate\\$k=6$} & \shortstack{HPDC-MaOEA\\$k=30$} \\
\midrule
DTLZ1 & \shortstack{$6.7408e+1^{+}$\\$(1.93e+1)$} & \shortstack{$1.0194e+2^{-}$\\$(2.14e+1)$} & \shortstack{$\mathbf{5.7687e+1}^{+}$\\$(1.80e+1)$} & \shortstack{$8.1313e+1^{=}$\\$(2.10e+1)$} & \shortstack{$6.0747e+1^{+}$\\$(1.62e+1)$} & \shortstack{$7.9755e+1$\\$(1.09e+1)$} \\
DTLZ2 & \shortstack{$1.1173e+0^{=}$\\$(5.45e-2)$} & \shortstack{$1.1484e+0^{=}$\\$(4.27e-2)$} & \shortstack{$1.1203e+0^{=}$\\$(5.61e-2)$} & \shortstack{$1.1979e+0^{-}$\\$(3.90e-2)$} & \shortstack{$\mathbf{1.0857e+0}^{+}$\\$(7.42e-2)$} & \shortstack{$1.1350e+0$\\$(4.89e-2)$} \\
DTLZ3 & \shortstack{$2.6646e+2^{=}$\\$(7.13e+1)$} & \shortstack{$3.7834e+2^{-}$\\$(8.48e+1)$} & \shortstack{$\mathbf{2.2547e+2}^{+}$\\$(5.49e+1)$} & \shortstack{$3.2741e+2^{=}$\\$(5.90e+1)$} & \shortstack{$2.5467e+2^{=}$\\$(7.79e+1)$} & \shortstack{$2.8568e+2$\\$(6.84e+1)$} \\
DTLZ4 & \shortstack{$1.0368e+0^{=}$\\$(4.26e-2)$} & \shortstack{$1.0608e+0^{=}$\\$(3.77e-2)$} & \shortstack{$1.0592e+0^{=}$\\$(4.75e-2)$} & \shortstack{$\mathbf{1.0199e+0}^{+}$\\$(4.38e-2)$} & \shortstack{$1.0634e+0^{=}$\\$(5.86e-2)$} & \shortstack{$1.0578e+0$\\$(5.34e-2)$} \\
DTLZ5 & \shortstack{$2.7068e-1^{=}$\\$(5.25e-2)$} & \shortstack{$3.3109e-1^{=}$\\$(4.78e-2)$} & \shortstack{$\mathbf{2.0792e-1}^{+}$\\$(5.04e-2)$} & \shortstack{$2.6667e-1^{=}$\\$(4.58e-2)$} & \shortstack{$2.8410e-1^{=}$\\$(4.78e-2)$} & \shortstack{$3.0345e-1$\\$(6.65e-2)$} \\
DTLZ6 & \shortstack{$7.3417e+0^{=}$\\$(6.33e-1)$} & \shortstack{$7.5855e+0^{-}$\\$(4.73e-1)$} & \shortstack{$\mathbf{6.3369e+0}^{=}$\\$(8.09e-1)$} & \shortstack{$6.5356e+0^{=}$\\$(7.58e-1)$} & \shortstack{$7.4829e+0^{-}$\\$(6.41e-1)$} & \shortstack{$6.6546e+0$\\$(9.39e-1)$} \\
DTLZ7 & \shortstack{$3.8239e+1^{-}$\\$(7.37e+0)$} & \shortstack{$4.2108e+1^{-}$\\$(7.08e+0)$} & \shortstack{$2.9216e+1^{-}$\\$(4.21e+0)$} & \shortstack{$2.5637e+1^{-}$\\$(5.00e+0)$} & \shortstack{$2.2735e+1^{-}$\\$(3.18e+0)$} & \shortstack{$\mathbf{1.8154e+1}$\\$(2.61e+0)$} \\
WFG1 & \shortstack{$5.6202e+0^{=}$\\$(8.82e-2)$} & \shortstack{$5.6481e+0^{-}$\\$(4.89e-2)$} & \shortstack{$5.5981e+0^{=}$\\$(1.05e-1)$} & \shortstack{$5.6149e+0^{=}$\\$(7.80e-2)$} & \shortstack{$\mathbf{5.5705e+0}^{=}$\\$(8.88e-2)$} & \shortstack{$5.6075e+0$\\$(7.93e-2)$} \\
WFG2 & \shortstack{$7.9480e+0^{-}$\\$(1.54e+0)$} & \shortstack{$7.0318e+0^{-}$\\$(1.58e+0)$} & \shortstack{$7.2512e+0^{-}$\\$(1.65e+0)$} & \shortstack{$\mathbf{5.1016e+0}^{=}$\\$(6.88e-1)$} & \shortstack{$6.8209e+0^{-}$\\$(2.06e+0)$} & \shortstack{$5.4393e+0$\\$(1.31e+0)$} \\
WFG3 & \shortstack{$\mathbf{2.2519e+0}^{+}$\\$(2.32e-1)$} & \shortstack{$2.3254e+0^{+}$\\$(1.85e-1)$} & \shortstack{$2.5155e+0^{+}$\\$(1.99e-1)$} & \shortstack{$2.6644e+0^{=}$\\$(1.17e-1)$} & \shortstack{$2.4822e+0^{+}$\\$(1.70e-1)$} & \shortstack{$2.7306e+0$\\$(1.52e-1)$} \\
WFG4 & -- & -- & -- & -- & -- & -- \\
WFG5 & \shortstack{$2.5608e+1^{-}$\\$(1.19e+0)$} & \shortstack{$2.3181e+1^{=}$\\$(1.06e+0)$} & \shortstack{$2.5693e+1^{-}$\\$(9.94e-1)$} & \shortstack{$2.3849e+1^{=}$\\$(1.38e+0)$} & \shortstack{$2.4357e+1^{-}$\\$(9.19e-1)$} & \shortstack{$\mathbf{2.2849e+1}$\\$(1.62e+0)$} \\
WFG6 & \shortstack{$2.4041e+1^{-}$\\$(1.66e+0)$} & \shortstack{$2.1493e+1^{-}$\\$(7.92e-1)$} & \shortstack{$2.3492e+1^{-}$\\$(1.36e+0)$} & \shortstack{$1.9109e+1^{=}$\\$(1.02e+0)$} & \shortstack{$2.0834e+1^{-}$\\$(9.11e-1)$} & \shortstack{$\mathbf{1.9097e+1}$\\$(6.18e-1)$} \\
WFG7 & \shortstack{$2.4322e+1^{-}$\\$(1.10e+0)$} & \shortstack{$2.2442e+1^{-}$\\$(1.43e+0)$} & \shortstack{$2.3916e+1^{-}$\\$(2.00e+0)$} & \shortstack{$\mathbf{2.0450e+1}^{=}$\\$(8.73e-1)$} & \shortstack{$2.2471e+1^{-}$\\$(2.16e+0)$} & \shortstack{$2.0592e+1$\\$(7.61e-1)$} \\
WFG8 & \shortstack{$2.4953e+1^{-}$\\$(1.02e+0)$} & \shortstack{$2.2718e+1^{-}$\\$(1.02e+0)$} & \shortstack{$2.4403e+1^{-}$\\$(1.29e+0)$} & \shortstack{$2.1098e+1^{=}$\\$(5.56e-1)$} & \shortstack{$2.2424e+1^{-}$\\$(1.08e+0)$} & \shortstack{$\mathbf{2.0738e+1}$\\$(8.26e-1)$} \\
WFG9 & \shortstack{$2.6155e+1^{-}$\\$(1.38e+0)$} & \shortstack{$2.4590e+1^{-}$\\$(1.41e+0)$} & \shortstack{$2.5690e+1^{-}$\\$(1.26e+0)$} & \shortstack{$2.4763e+1^{-}$\\$(1.35e+0)$} & \shortstack{$2.4951e+1^{-}$\\$(1.72e+0)$} & \shortstack{$\mathbf{2.3053e+1}$\\$(1.30e+0)$} \\
\midrule
W/T/L & 7/6/2 & 10/4/1 & 7/4/4 & 3/11/1 & 8/4/3 & -- \\
Best mean & 1 & 0 & 4 & 3 & 2 & 5 \\
\bottomrule
\end{tabular}
\par\smallskip
\begin{minipage}{\linewidth}
\scriptsize Cells report mean IGD (standard deviation). Symbols and W/T/L follow the main tables. HPC denotes hybrid PBI grouping. The enlarged-$k$ HPC-only configuration isolates addition of candidate selection when compared with the full method. WFG4 is unfilled because the source row records $M=10$; all counts use the remaining 15 problems.
\end{minipage}
\end{table*}

\TODO{Insert the candidate-only configuration at $k=15,30$ and complete
the matched component-removal comparisons. Leave unverified reduced-$k$
full-method entries blank until their implementation identity is resolved.
The joint interaction of the two modules remains to be tested.}

\subsection{Analysis of Hybrid PBI-Based Quality Grouping}
\label{sec:exp:hpc}

\paragraph{Positive-group composition.}
Hybrid PBI grouping selected a positive group distinct from either
single-view ranking. This was examined using the archived snapshot study
on DTLZ2, DTLZ4, DTLZ7, WFG3 and WFG7 at $M=10,20$, with five runs for
each of two search behaviours: Hybrid and the native representative-based
binary grouping. The study comprised 100 trajectories and 6,700
checkpoints. The same snapshots were used for offline comparisons of
positive-group definitions, each retaining the top 25\% except for the
native binary group.

Along the 50 Hybrid trajectories, mean Jaccard overlap was 0.5470 between
the Hybrid group and the direction-score group, and 0.3996 between the
Hybrid group and the representative-margin group. Along the 50
native-group trajectories, the corresponding values were 0.5699 and
0.4244. These run-averaged overlaps establish a change in group
composition under both recorded search behaviours. Here, the continuous
representative margin is an equal-quota comparator. The production
Hybrid score combines the direction score with the binary label, as
specified in Section~\ref{sec:method}; it does not combine two continuous
scores.

\paragraph{Association with subsequent retention.}
The completed good-group precision study evaluated the same five
problems at $M=10,20$ with 25 runs per configuration, giving 250
trajectories at $FE_{\max}=500$. At each checkpoint, the Hybrid,
direction-score and representative-margin rankings selected the top
25\% of the current population. Precision was the fraction of selected
solutions that remained in the final population of the recorded Hybrid
trajectory. Checkpoint precision was averaged within each run and one
of four $FE/FE_{\max}$ stages, followed by equal averaging across runs
and configurations.

\begin{table}[t]
\centering
\caption{Precision@25\% for final-population retention.}
\label{tab:exp:precision}
\footnotesize
\setlength{\tabcolsep}{3pt}
\renewcommand{\arraystretch}{1.12}
\begin{tabular}{lccc}
\toprule
$FE/FE_{\max}$ & Hybrid & Direction & Margin \\
\midrule
$[0,0.25]$ & 0.08547 & 0.05949 & 0.07949 \\
$(0.25,0.50]$ & 0.11804 & 0.10751 & 0.12414 \\
$(0.50,0.75]$ & 0.29011 & 0.28834 & 0.27126 \\
$(0.75,1.00]$ & 0.65431 & 0.65458 & 0.62980 \\
\midrule
All stages & 0.28698 & 0.27748 & 0.27617 \\
\bottomrule
\end{tabular}
\par\smallskip
\begin{minipage}{\linewidth}
\scriptsize Each stage averages 250 run-level values per view. The last row gives equal weight to the four stages. Larger values mean a greater fraction of selected solutions remained in the final population.
\end{minipage}
\end{table}

Hybrid achieved mean Precision@25\% of 0.28698, compared with 0.27748
for the direction score and 0.27617 for the representative margin
(Table~\ref{tab:exp:precision}). Its advantage was not uniform across
stages: the representative margin had the highest mean in the second
stage, while the direction score and Hybrid were nearly equal in the
last stage. The native binary group selected 49.05\% of the population
on average and had precision 0.31905; its different selection fraction
requires it to be reported separately from the top-25\% comparisons.

The two single-view groups also contained distinct future-retained
solutions. Across 40 problem--objective--stage cells, 33 passed the
Holm-adjusted test requiring both views to contribute unique
future-retained solutions in a majority of runs. The stricter test
requiring Hybrid precision to exceed both single-view precisions passed
in one cell. Thus, the evidence supports nonredundant information in the
two views and a modest aggregate retention advantage for Hybrid along
the recorded trajectories. The endpoint measured here is future
population membership, rather than final algorithm-level IGD.

\paragraph{Effect of the mixing schedule.}
The offline snapshot comparison additionally evaluated a fixed
equal-weight mixture and a reversed schedule. Along the Hybrid
trajectories, their mean positive-group overlaps with the original
schedule were 0.99987 and 0.98806, respectively. In these snapshots,
group membership was therefore largely preserved by the alternative
schedules. These observations complement the ordering rationale in the
method section; they do not establish an endpoint advantage for the
time-varying schedule.

\TODO{Add matched closed-loop comparisons of the grouping schedules and
independent utility results after the reference-stability check is resolved.}

\subsection{Analysis of Dual-Mode Candidate Selection}
\label{sec:exp:candidate}

\paragraph{Batch expenditure and selection criteria.}
The candidate-value probe compared five policies on DTLZ2, DTLZ7, WFG3
and WFG7 with $M=20$, $N=100$ and $FE_{\max}=300$. Each policy had ten
runs per problem, paired by the problem and run seed. V0 applied the
REMO selection rule within the probe host. V1 selected the six highest
relation scores from the accumulated pool. V2 used exploration-only
selection, V3 indicator-only selection, and V4 the full dual-mode
strategy. V0 is a transplanted selection rule, not the original REMO
algorithm. V1--V4 used the same accumulated-pool construction and a
six-candidate batch bound; their search trajectories generated their
own pools.

All 200 runs completed at 300 evaluations. V1--V4 selected exactly six
candidates in every non-truncated iteration, with the final batch
reduced to two by the remaining budget. Their mean batch size was
5.882. Under V0, the batch size ranged from 1 to 38 and averaged 10.386.
This comparison directly demonstrates the expenditure control supplied
by the explicit batch bound.

\paragraph{Batch properties and candidate utility.}
Batch spread was the mean pairwise Euclidean distance between selected
candidates after scaling each decision variable by its bounds. Retention was the fraction of evaluated candidates accepted
by the host's population-selection operation. The within-pool gain
ratio compared the selected batch's reduction in mean reference-to-set
distance with that of a greedy reference batch. This diagnostic used
300 objective-space reference points and a 400-candidate subset that
included all selected candidates. Its objective evaluations were used
only for offline assessment and were not returned to the search.
Retention and gain ratios below use iterations with
$FE/FE_{\max}\geq0.5$; batch spread and size use the full trajectory.

\begin{table*}[t]
\centering
\caption{Candidate-value probe at $M=20$: final IGD and batch properties.}
\label{tab:exp:candidate}
\scriptsize
\setlength{\tabcolsep}{4pt}
\renewcommand{\arraystretch}{1.12}
\begin{tabular}{lccccc}
\toprule
Problem / metric & V0: REMO rule & V1: relation top-six & V2: exploration & V3: indicator & V4: dual mode \\
\midrule
DTLZ2 & \shortstack{$1.2052e+00$\\$(6.00e-02)$} & \shortstack{$1.2185e+00$\\$(3.75e-02)$} & \shortstack{$1.2165e+00$\\$(4.32e-02)$} & \shortstack{$1.1124e+00$\\$(5.25e-02)$} & \shortstack{$\mathbf{1.1109e+00}$\\$(4.67e-02)$} \\
DTLZ7 & \shortstack{$2.6138e+01$\\$(4.68e+00)$} & \shortstack{$\mathbf{1.7364e+01}$\\$(5.53e+00)$} & \shortstack{$2.0498e+01$\\$(9.13e+00)$} & \shortstack{$3.2509e+01$\\$(6.80e+00)$} & \shortstack{$1.9675e+01$\\$(2.16e+00)$} \\
WFG3 & \shortstack{$2.7212e+00$\\$(1.18e-01)$} & \shortstack{$2.8720e+00$\\$(8.43e-02)$} & \shortstack{$2.8239e+00$\\$(1.10e-01)$} & \shortstack{$\mathbf{2.6946e+00}$\\$(1.63e-01)$} & \shortstack{$2.7877e+00$\\$(1.40e-01)$} \\
WFG7 & \shortstack{$2.0776e+01$\\$(6.51e-01)$} & \shortstack{$2.1052e+01$\\$(6.09e-01)$} & \shortstack{$2.1241e+01$\\$(1.04e+00)$} & \shortstack{$2.2036e+01$\\$(1.18e+00)$} & \shortstack{$\mathbf{2.0318e+01}$\\$(1.42e+00)$} \\
\midrule
Mean batch size & 10.386 & 5.882 & 5.882 & 5.882 & 5.882 \\
Mean batch spread & 1.806 & 1.637 & 2.031 & 1.031 & 1.479 \\
Late retention & 0.756 & 0.706 & 0.676 & 0.761 & 0.770 \\
Late gain ratio & 0.123 & 0.086 & 0.081 & 0.059 & 0.105 \\
\bottomrule
\end{tabular}
\par\smallskip
\begin{minipage}{\linewidth}
\scriptsize IGD entries are means (sample standard deviations) over ten runs per problem and policy. Batch statistics are equally averaged across the 40 runs of each policy. Only retention and gain ratio use the late-stage restriction. The V0 rule uses the probe host, and the greedy-reference diagnostic is restricted to the sampled pool.
\end{minipage}
\end{table*}

The two modes produced different batch properties. Exploration-only
selection had the largest mean spread (2.031), whereas indicator-only
selection had the smallest (1.031). V4 attained late-stage retention
of 0.770 and a gain ratio of 0.105, compared with 0.706 and 0.086 for
relation top-six selection. V4's mean spread was 1.479, between the two
single-mode values. These measurements describe the selected batches
and are interpreted jointly with the problemwise endpoint IGD in
Table~\ref{tab:exp:candidate}.

V4 had lower mean final IGD than V1 on DTLZ2, WFG3 and WFG7, but higher
mean IGD on DTLZ7 (19.675 versus 17.364). Indicator-only selection
illustrates why retention alone is insufficient: its late-stage
retention was 0.761, but its mean IGD on DTLZ7 was 32.509. The two
criteria therefore differ in both batch composition and final quality,
and the benefit of their combination depends on the problem.

\paragraph{Within-host routing comparison.}
A separate routing experiment compared exploration-only,
indicator-only, linearly scheduled switching and equal-probability
switching on all 16 problems at $M=10,20$. These results used the
\texttt{UniformMix} implementation of the routing-study host; they are
reported separately from the \texttt{UniformMix\_Original} main results.
Table~\ref{tab:exp:routing} retains the IGD and available IGD$^+$
comparisons within that host.

Equal-probability switching obtained reported IGD W/T/L counts of
5/9/2 against either single mode at $M=10$. At $M=20$, the counts were
4/11/1 against exploration-only and 5/11/0 against indicator-only.
The DTLZ7 mean IGD was 18.615 under equal-probability switching and
33.514 under indicator-only selection. Against linear switching, the
counts were 0/16/0 at $M=10$ and 1/15/0 at $M=20$. Thus, these results
support using both criteria within this host, while showing little
separation between the two switching schedules.

\begin{table}[t]
\centering
\caption{Equal-probability switching versus alternative routing policies in the routing-study host.}
\label{tab:exp:routing}
\scriptsize
\setlength{\tabcolsep}{3pt}
\renewcommand{\arraystretch}{1.12}
\begin{tabular}{lccc}
\toprule
$M$ & Indicator & Comparator & W/T/L \\
\midrule
10 & IGD & Exploration-only & 5/9/2 \\
10 & IGD & Indicator-only & 5/9/2 \\
10 & IGD & Linear switching & 0/16/0 \\
\midrule
20 & IGD & Exploration-only & 4/11/1 \\
20 & IGD & Indicator-only & 5/11/0 \\
20 & IGD & Linear switching & 1/15/0 \\
\midrule
20 & IGD$^+$ & Exploration-only & 3/12/1 \\
20 & IGD$^+$ & Indicator-only & 5/9/2 \\
20 & IGD$^+$ & Linear switching & 1/15/0 \\
\bottomrule
\end{tabular}
\par\smallskip
\begin{minipage}{\linewidth}
\scriptsize Counts use all 16 problems and the exported annotations relative to UniformMix. These are separate host-specific comparisons; they are not ablations of the UniformMix\_Original entries in the main tables.
\end{minipage}
\end{table}

\TODO{Complete the relation-top-six and single-mode ablations for the
final main-comparison implementation, together with the missing
10-objective candidate-value probe.}

\subsection{Robustness to Core Controls}
\label{sec:exp:sensitivity}

\paragraph{Stability of positive-group membership.}
The offline grouping study tested removal of 5\% and 10\% of the
population before reconstructing the group. Mean retained-set Jaccard
overlap for Hybrid was 0.97074 and 0.95471 along the Hybrid trajectories,
respectively. The native-group trajectories yielded 0.97081 and
0.95470. Group membership was therefore stable under these local
population perturbations. This test concerns label stability at fixed
algorithm controls.

\paragraph{Additional evaluation budget.}
The available 500-evaluation comparison covered the 16 10-objective
problems with the \texttt{UniformMix} host. Its reported W/T/L counts
were 10/4/2 against REMO, 7/4/5 against PIEA, 11/1/4 against MCEAD,
7/4/5 against R2AEA, 8/4/4 against HES-EA and 15/1/0 against PCSAEA.
This provides additional within-budget performance evidence. A
matched budget comparison for the final implementation remains open.

\TODO{Add final-implementation sensitivity results for the core
grouping and candidate-allocation controls, matched FE-budget comparisons,
and the wall-clock shares of the two modules.}

\clearpage

\section{Conclusion}
\label{sec:conclusion}
\TODO{Reinforce the two design contributions: hybrid PBI-based grouping combines
fine-grained ordering with representative-based group priority, and dual-mode
candidate selection offers exploration and indicator criteria within an
explicit evaluation-batch bound. State the strongest verified algorithm-level
result and its benchmark scope. Close by connecting the group-composition and
candidate-batch findings to these two design roles, keeping empirical performance
claims at the level supported by their matched comparisons.}

\begin{thebibliography}{9}

\bibitem{ref:moead}
Q.~Zhang and H.~Li,
``MOEA/D: A multiobjective evolutionary algorithm based on decomposition,''
\emph{IEEE Trans. Evol. Comput.}, vol.~11, no.~6, pp.~712--731, 2007.

\bibitem{ref:remo}
\TODO{REMO -- relation-model-based expensive multi-objective optimization.}

\bibitem{ref:rsea}
\TODO{RSEA -- radar-grid based selection evolutionary algorithm.}

\bibitem{ref:piea}
Y.~Li, W.~Li, S.~Li, and Y.~Zhao,
``PIEA,'' \emph{Information Sciences}, 2024.
\TODO{complete volume and pages.}

\bibitem{ref:sde}
M.~Li, S.~Yang, and X.~Liu,
``Shift-based density estimation for Pareto-based algorithms in
many-objective optimization,''
\emph{IEEE Trans. Evol. Comput.}, vol.~18, no.~3, pp.~348--365, 2014.

\bibitem{ref:pcsaea}
\TODO{PC-SAEA -- pairwise-comparison-based surrogate-assisted EA,
Swarm Evol. Comput., 2023, doi:10.1016/j.swevo.2023.101323.}

\bibitem{ref:r2aea}
\TODO{R2AEA -- two-stage relation-based expensive optimization.}

\bibitem{ref:platemo}
Y.~Tian, R.~Cheng, X.~Zhang, and Y.~Jin,
``PlatEMO: A MATLAB platform for evolutionary multi-objective optimization,''
\emph{IEEE Comput. Intell. Mag.}, vol.~12, no.~4, pp.~73--87, 2017.

\end{thebibliography}

\end{document}
